\section*{Practice 2--3: Relations}

%% ============================================================
%%  1. Relations
%% ============================================================

\subsection*{Problem 1.1}

\textbf{(a)} Divisors of 12 in $\mathbb{N}$:
\[
  \boxed{\{(1,12),(2,12),(3,12),(4,12),(6,12),(12,12)\}}
\]

\textbf{(b)} Pairs $(d,n)$ with $d,n\in\{2,3,4,5,6\}$ and $d\mid n$:
\[
  \boxed{\{(2,2),(2,4),(2,6),(3,3),(3,6),(4,4),(5,5),(6,6)\}}
\]

\textbf{(c)} Triples $(x,y,z)$ with $x=y+z$, all in $\{1,2,3\}$:
\[
  \boxed{\{(2,1,1),(3,1,2),(3,2,1)\}}
\]

%% ────────────────────────────────────────

\subsection*{Problem 1.2}

\textbf{(a)} $\boxed{D=\{a,c,d\},\quad R=\{1,2,4,5\}}$.

\textbf{(b)} Pattern $(n,2n)$: $\boxed{D=\mathbb{N},\quad R=2\mathbb{N}}$.

\textbf{(c)} $x=y^2$: $\boxed{D=[0,+\infty),\quad R=\mathbb{R}}$.

\textbf{(d)} $x^2+y^2\le 16$: $\boxed{D=[-4,4],\quad R=[-4,4]}$.

\textbf{(e)} $(x,y)\in\mathbb{N}^2,\; x\le y$: $\boxed{D=\mathbb{N},\quad R=\mathbb{N}}$.

\textbf{(f)} $(x,y)\in\mathbb{N}^2,\; y\le x$: $\boxed{D=\mathbb{N},\quad R=\mathbb{N}}$.

\textbf{(g)} $x+y\ge 0$ on $\mathbb{R}^2$: $\boxed{D=\mathbb{R},\quad R=\mathbb{R}}$.

\textbf{(h)} $2x\ge 3y$ on $\mathbb{R}^2$: $\boxed{D=\mathbb{R},\quad R=\mathbb{R}}$.

%% ============================================================
%%  2. Representation of Relations
%% ============================================================

\subsection*{Problem 2.1}

On $A=\{a,b,c,d,e\}$.

\textbf{(a)} Symmetric; undirected edges: $a$--$b$, $a$--$c$, $b$--$c$, $d$--$e$.
\[
  M = \begin{pmatrix}
  0&1&1&0&0\\1&0&1&0&0\\1&1&0&0&0\\0&0&0&0&1\\0&0&0&1&0
  \end{pmatrix}
\]

\textbf{(b)} Undirected edges: $a$--$b$, $a$--$c$, $b$--$c$, $c$--$d$. Vertex $e$ isolated.
\[
  M = \begin{pmatrix}
  0&1&1&0&0\\1&0&1&0&0\\1&1&0&1&0\\0&0&1&0&0\\0&0&0&0&0
  \end{pmatrix}
\]

\textbf{(c)} Identical to part (b); same graph and matrix.

%% ────────────────────────────────────────

\subsection*{Problem 2.2}

On $A=\{a,b,c,d,e,f\}$.  Rows/columns in order $a,b,c,d,e,f$.

\textbf{(a)} Directed edges: $a\!\to\!b,\; b\!\to\!c,\; a\!\to\!c,\; b\!\to\!e,\; c\!\to\!f,\; c\!\to\!d,\; d\!\to\!f,\; f\!\to\!e$.
\[
  M = \begin{pmatrix}
  0&1&1&0&0&0\\0&0&1&0&1&0\\0&0&0&1&0&1\\0&0&0&0&0&1\\0&0&0&0&0&0\\0&0&0&0&1&0
  \end{pmatrix}
\]

\textbf{(b)} Directed edges: $a\!\to\!b,\; b\!\to\!c,\; d\!\to\!c,\; d\!\to\!e,\; f\!\to\!e,\; f\!\to\!a,\; b\!\to\!e$.
\[
  M = \begin{pmatrix}
  0&1&0&0&0&0\\0&0&1&0&1&0\\0&0&0&0&0&0\\0&0&1&0&1&0\\0&0&0&0&0&0\\1&0&0&0&1&0
  \end{pmatrix}
\]

\textbf{(c)} Directed edges: $a\!\to\!c,\; c\!\to\!c,\; c\!\to\!d,\; c\!\to\!e,\; e\!\to\!c,\; e\!\to\!e,\; e\!\to\!f,\; a\!\to\!e$. Vertex $b$ isolated.
\[
  M = \begin{pmatrix}
  0&0&1&0&1&0\\0&0&0&0&0&0\\0&0&1&1&1&0\\0&0&0&0&0&0\\0&0&1&0&1&1\\0&0&0&0&0&0
  \end{pmatrix}
\]

%% ============================================================
%%  3. Operations on Relations
%% ============================================================

\subsection*{Problem 3.1}

\emph{Proof sketch} (standard properties of ${}^{-1}$ and $\circ$):

(i) $(\rho^{-1})^{-1}=\rho$: $(x,y)\in(\rho^{-1})^{-1}\Leftrightarrow (y,x)\in\rho^{-1}\Leftrightarrow (x,y)\in\rho$.

(ii) $(\rho\circ\sigma)^{-1}=\sigma^{-1}\circ\rho^{-1}$: $(z,x)\in(\rho\circ\sigma)^{-1} \Leftrightarrow \exists y\!:\,(x,y)\in\rho\land(y,z)\in\sigma \Leftrightarrow \exists y\!:\,(z,y)\in\sigma^{-1}\land(y,x)\in\rho^{-1} \Leftrightarrow (z,x)\in\sigma^{-1}\circ\rho^{-1}$. \qed

%% ────────────────────────────────────────

\subsection*{Problem 3.2}

$\rho=\{(x,y)\mid y=x^2+5\}$, $\sigma=\{(x,y)\mid y=3x\}$ on $\mathbb{R}$.

$\boxed{\rho^{-1}=\{(x,y)\mid x=y^2+5\}}$ (i.e.\ $y=\pm\sqrt{x-5}$, defined for $x\ge 5$).

$\boxed{\sigma^{-1}=\{(x,y)\mid y=x/3\}}$.

$\rho\circ\sigma$: $x\xrightarrow{\rho}x^2+5\xrightarrow{\sigma}3(x^2+5)$. \;$\boxed{\rho\circ\sigma=\{(x,\,3x^2+15)\}}$.

$\sigma\circ\rho$: $x\xrightarrow{\sigma}3x\xrightarrow{\rho}9x^2+5$. \;$\boxed{\sigma\circ\rho=\{(x,\,9x^2+5)\}}$.

%% ────────────────────────────────────────

\subsection*{Problem 3.3}

$\rho=\{(x,x+2)\}$, $\sigma=\{(x,x-2)\}$ on $\mathbb{Z}$.

$\rho\circ\sigma$: $x\xrightarrow{\rho}x+2\xrightarrow{\sigma}x$. \;$\boxed{\rho\circ\sigma=\Delta_{\mathbb{Z}}}$ (identity).

$\sigma\circ\rho$: $x\xrightarrow{\sigma}x-2\xrightarrow{\rho}x$. \;$\boxed{\sigma\circ\rho=\Delta_{\mathbb{Z}}}$.

%% ────────────────────────────────────────

\subsection*{Problem 3.4}

$\rho=\{(x,y)\in\mathbb{Z}^2\mid x+y \text{ is odd}\}$.

For $\rho^2$: need $z$ with $x+z$ odd and $z+y$ odd, so $x\equiv y\pmod{2}$.
$\boxed{\rho^2=\{(x,y)\in\mathbb{Z}^2\mid x+y \text{ is even}\}}$.

$\boxed{\rho^3=\rho}$, and in general $\rho^{2k}=\rho^2$, $\rho^{2k+1}=\rho$.

%% ────────────────────────────────────────

\subsection*{Problem 3.5}

$\rho=\{(x,x+1)\mid x\in\mathbb{N}\}$.

\[\boxed{\rho^k=\{(x,\,x+k)\mid x\in\mathbb{N}\},\quad k\ge 1.}\]

%% ────────────────────────────────────────

\subsection*{Problem 3.6}

$\rho=\{(1,2),(2,3),(3,4),(4,5)\}$.

$\boxed{\rho^2=\{(1,3),(2,4),(3,5)\}}$.

$\boxed{\rho^3=\{(1,4),(2,5)\}}$.

$\boxed{\rho^4=\{(1,5)\}}$.

$\boxed{\rho^5=\varnothing}$.

%% ============================================================
%%  4. Properties of Relations
%% ============================================================

\subsection*{Problem 4.1}

\begin{tabular}{l|cccc}
& Refl. & Symm. & Antisymm. & Trans. \\\hline
(a) $x^2+y^2=1$, $\mathbb{R}$ & No & Yes & No & No \\
(b) $x^2=y^2$, $\mathbb{R}$   & Yes & Yes & No & Yes \\
(c) $x\le y$, $\mathbb{Z}$    & Yes & No  & Yes & Yes \\
(d) $x\le y$, $\mathbb{N}$    & Yes & No  & Yes & Yes \\
(e) $\gcd=1$, $\mathbb{N}$    & No  & Yes & No  & No \\
(f) $\gcd=1$, $\mathbb{Z}$    & No  & Yes & No  & No
\end{tabular}

\smallskip
(b) is an equivalence relation; (c) and (d) are linear (total) orders.

%% ────────────────────────────────────────

\subsection*{Problem 4.2}

$\rho=\{(x,y)\in\mathbb{Z}^2\mid |x-y|\text{ is even}\}$.

Reflexive -- Yes; Symmetric -- Yes; Antisymmetric -- No; Transitive -- Yes.

$\boxed{\text{Equivalence relation. Classes: even integers and odd integers.}}$

%% ────────────────────────────────────────

\subsection*{Problem 4.3}

Preservation of \textbf{reflexivity} ($R,S$ reflexive on $A$):

(a) $R\cap S$: \textbf{True.} \quad
(b) $R\cup S$: \textbf{True.} \quad
(c) $R\circ S$: \textbf{True.}

(d) $R\setminus S$: \textbf{False} ($(x,x)\in R\cap S$ so $(x,x)\notin R\setminus S$). \quad
(e) $R\oplus S$: \textbf{False} (same reason).

%% ────────────────────────────────────────

\subsection*{Problem 4.4}

Preservation of \textbf{symmetry} ($R,S$ symmetric):

(a) $R\cap S$: \textbf{True.} \quad
(b) $R\cup S$: \textbf{True.}

(c) $R\circ S$: \textbf{False.} ($R=\{(1,2),(2,1)\}$, $S=\{(2,3),(3,2)\}$: $(1,3)\in R\circ S$ but $(3,1)\notin R\circ S$.)

(d) $R\setminus S$: \textbf{True.} \quad
(e) $R\oplus S$: \textbf{True.}

%% ────────────────────────────────────────

\subsection*{Problem 4.5}

Preservation of \textbf{antisymmetry} ($R,S$ antisymmetric):

(a) $R\cap S$: \textbf{True.} \quad
(b) $R\cup S$: \textbf{False} ($R=\{(1,2)\}$, $S=\{(2,1)\}$).

(c) $R\circ S$: \textbf{False.} \quad
(d) $R\setminus S$: \textbf{True} ($R\setminus S\subseteq R$).

(e) $R\oplus S$: \textbf{False} (same counterexample as (b)).

%% ────────────────────────────────────────

\subsection*{Problem 4.6}

Preservation of \textbf{transitivity} ($R,S$ transitive):

(a) $R\cap S$: \textbf{True.}

(b) $R\cup S$: \textbf{False} ($R=\{(1,2)\}$, $S=\{(2,3)\}$; $(1,3)\notin R\cup S$).

(c) $R\circ S$: \textbf{False.}

(d) $R\setminus S$: \textbf{False} ($R=\{(1,2),(2,3),(1,3)\}$, $S=\{(1,3)\}$; $R\setminus S=\{(1,2),(2,3)\}$, not transitive).

(e) $R\oplus S$: \textbf{False.}

%% ============================================================
%%  5. Order Relations
%% ============================================================

\subsection*{Problem 5.1}

Construct a poset with a unique minimal element but no least element.

\emph{Note:} In any \emph{finite} poset a unique minimal element is automatically a least element, so we need an infinite poset.

$\boxed{\text{Example: }A=\{m\}\cup\{a_1,a_2,a_3,\ldots\},\text{ with }a_1<a_2<a_3<\cdots}$
and $m$ incomparable to every $a_i$. Every $a_i$ has $a_{i-1}<a_i$, so no $a_i$ is minimal. Hence $m$ is the unique minimal element. But $m\not\le a_i$ for any $i$, so $m$ is not a least element.

%% ────────────────────────────────────────

\subsection*{Problem 5.2}

\emph{Proof sketch.} Let $\rho$ be a partial order on $A$. Show $\rho^{-1}$ is reflexive, antisymmetric, transitive.
Reflexive: $(x,x)\in\rho\Rightarrow(x,x)\in\rho^{-1}$.
Antisymmetric: $(x,y),(y,x)\in\rho^{-1}\Rightarrow(y,x),(x,y)\in\rho\Rightarrow x=y$.
Transitive: $(x,y),(y,z)\in\rho^{-1}\Rightarrow(z,y),(y,x)\in\rho\Rightarrow(z,x)\in\rho\Rightarrow(x,z)\in\rho^{-1}$. \qed

%% ────────────────────────────────────────

\subsection*{Problem 5.3}

\emph{Proof sketch.} Let $\{\rho_i\}$ be partial orders on $A$, $\rho=\bigcap_i\rho_i$.
Reflexive: $(x,x)\in\rho_i$ for all $i$.
Antisymmetric: $(x,y),(y,x)\in\rho\subseteq\rho_i\Rightarrow x=y$.
Transitive: $(x,y),(y,z)\in\rho\subseteq\rho_i\Rightarrow(x,z)\in\rho_i$ for all $i\Rightarrow(x,z)\in\rho$. \qed

%% ────────────────────────────────────────

\subsection*{Problem 5.4}

\emph{Proof sketch.} $\rho_A=\rho\cap(A\times A)$ inherits reflexivity (for $x\in A$), antisymmetry, and transitivity from $\rho$ since $\rho_A\subseteq\rho$ and all elements in consideration belong to $A$. \qed

%% ────────────────────────────────────────

\subsection*{Problem 5.5}

$\boxed{\rho_1\cup\rho_2\text{ is a linear order iff }\rho_1\subseteq\rho_2\text{ or }\rho_2\subseteq\rho_1.}$

If one contains the other, the union equals the larger one. Conversely, if neither contains the other, antisymmetry of the union can fail.

%% ────────────────────────────────────────

\subsection*{Problem 5.6}

\emph{Proof sketch.} Induction on $|A|=n$. Base $n=1$: trivial. Step: pick $a\in A$; by hypothesis $A\setminus\{a\}$ has a linear order $a_1<\cdots<a_n$. Insert $a$ at the end: $a_1<\cdots<a_n<a$. This is a linear order on $A$. \qed

%% ============================================================
%%  6. Equivalence Relations
%% ============================================================

\subsection*{Problem 6.1}

\textbf{(a)} $xy>0$ on $\mathbb{Q}$: \textbf{No.} Not reflexive ($0\cdot 0=0\not>0$).

\textbf{(b)} $xy\in\mathbb{Z}$ on $\mathbb{Q}$: \textbf{No.} Not transitive ($\tfrac{1}{6}\cdot 6\in\mathbb{Z}$, $6\cdot\tfrac{1}{2}\in\mathbb{Z}$, but $\tfrac{1}{6}\cdot\tfrac{1}{2}\notin\mathbb{Z}$).

\textbf{(c)} $(a\mid b)\lor(a>b)$ on $\mathbb{N}$: \textbf{No.} Not symmetric ($(6,4)\in\rho$ since $6>4$, but $(4,6)\notin\rho$ since $4\nmid 6$ and $4<6$).

\textbf{(d)} $|a-b|\le 2$ on $\mathbb{N}$: \textbf{No.} Not transitive ($|1-3|=2\le 2$, $|3-5|=2\le 2$, but $|1-5|=4>2$).

%% ────────────────────────────────────────

\subsection*{Problem 6.2}

(Equivalence relation $\Leftrightarrow$ partition.)

\emph{Proof sketch.}
$(\Rightarrow)$ Classes cover $A$ (reflexivity: $x\in[x]$). If $[a]\cap[b]\ne\varnothing$, pick $c$ in both; then $a\rho c$ and $c\rho b$, so $a\rho b$ and $[a]=[b]$.
$(\Leftarrow)$ Define $x\rho y$ iff $x,y$ are in the same block; reflexivity, symmetry, transitivity are immediate. \qed

%% ────────────────────────────────────────

\subsection*{Problem 6.3}

\textbf{$\mathbb{Z}_2$:}
\[
\begin{array}{c|cc}
+ & [0] & [1] \\\hline
{[0]} & [0] & [1] \\
{[1]} & [1] & [0]
\end{array}
\qquad
\begin{array}{c|cc}
\cdot & [0] & [1] \\\hline
{[0]} & [0] & [0] \\
{[1]} & [0] & [1]
\end{array}
\]

\textbf{$\mathbb{Z}_3$:}
\[
\begin{array}{c|ccc}
+ & [0] & [1] & [2] \\\hline
{[0]} & [0] & [1] & [2] \\
{[1]} & [1] & [2] & [0] \\
{[2]} & [2] & [0] & [1]
\end{array}
\qquad
\begin{array}{c|ccc}
\cdot & [0] & [1] & [2] \\\hline
{[0]} & [0] & [0] & [0] \\
{[1]} & [0] & [1] & [2] \\
{[2]} & [0] & [2] & [1]
\end{array}
\]

\textbf{$\mathbb{Z}_4$:}
\[
\begin{array}{c|cccc}
+ & [0] & [1] & [2] & [3] \\\hline
{[0]} & [0] & [1] & [2] & [3] \\
{[1]} & [1] & [2] & [3] & [0] \\
{[2]} & [2] & [3] & [0] & [1] \\
{[3]} & [3] & [0] & [1] & [2]
\end{array}
\qquad
\begin{array}{c|cccc}
\cdot & [0] & [1] & [2] & [3] \\\hline
{[0]} & [0] & [0] & [0] & [0] \\
{[1]} & [0] & [1] & [2] & [3] \\
{[2]} & [0] & [2] & [0] & [2] \\
{[3]} & [0] & [3] & [2] & [1]
\end{array}
\]

%% ────────────────────────────────────────

\subsection*{Problem 6.4}

$[a]\in\mathbb{Z}_n$ is invertible iff $\gcd(a,n)=1$.

\emph{Proof sketch.}
$(\Rightarrow)$ $ab\equiv 1\pmod{n}$ means $ab - kn = 1$, so $\gcd(a,n)=1$ by B\'ezout.
$(\Leftarrow)$ $\gcd(a,n)=1$ implies $\exists\,b,k$: $ab+kn=1$, so $ab\equiv 1\pmod{n}$. \qed

%% ============================================================
%%  7. Operations on Equivalence Relations (Closures)
%% ============================================================

\subsection*{Problem 7.1}

On $A=\{a,b,c\}$, $\Delta_A=\{(a,a),(b,b),(c,c)\}$.

\textbf{(a)} $\rho=\{(a,b),(b,a)\}$.
$r(\rho)=\rho\cup\Delta_A=\{(a,a),(b,b),(c,c),(a,b),(b,a)\}$.
$s(\rho)=\rho$ (already symmetric).
$t(\rho)=\{(a,b),(b,a),(a,a),(b,b)\}$.

\textbf{(b)} $\rho=\{(a,b),(b,c)\}$.
$r(\rho)=\{(a,a),(b,b),(c,c),(a,b),(b,c)\}$.
$s(\rho)=\{(a,b),(b,a),(b,c),(c,b)\}$.
$t(\rho)=\{(a,b),(b,c),(a,c)\}$.

\textbf{(c)} $\rho=\{(a,a),(a,b),(c,b),(c,a)\}$.
$r(\rho)=\rho\cup\{(b,b),(c,c)\}$.
$s(\rho)=\rho\cup\{(b,a),(b,c),(a,c)\}$.
$t(\rho)=\rho$ (already transitive: $(c,a)\circ(a,b)=(c,b)\in\rho$).

\textbf{(d)} $\rho=\{(a,a),(a,c),(b,c)\}$.
$r(\rho)=\rho\cup\{(b,b),(c,c)\}$.
$s(\rho)=\rho\cup\{(c,a),(c,b)\}$.
$t(\rho)=\rho$ (already transitive).

\textbf{(e)} $\rho=\{(a,b),(c,a),(b,c)\}$.
$r(\rho)=\rho\cup\Delta_A$.
$s(\rho)=\{(a,b),(b,a),(c,a),(a,c),(b,c),(c,b)\}$.
$t(\rho)=A\times A$ (the full cycle $a\to b\to c\to a$ generates all 9 pairs).

\textbf{(f)} $\rho=\{(a,b),(c,b),(b,c)\}$.
$r(\rho)=\rho\cup\Delta_A$.
$s(\rho)=\{(a,b),(b,a),(c,b),(b,c)\}$.
$t(\rho)=\{(a,b),(a,c),(b,c),(b,b),(c,b),(c,c)\}$.

%% ────────────────────────────────────────

\subsection*{Problem 7.2}

$A=\{1,2,3,4\}$, $\rho=\{(1,2),(3,1),(3,2),(2,4)\}$.

$\boxed{r(\rho)=\{(1,1),(2,2),(3,3),(4,4),(1,2),(2,4),(3,1),(3,2)\}}$.

$\boxed{s(\rho)=\{(1,2),(2,1),(1,3),(3,1),(2,3),(3,2),(2,4),(4,2)\}}$.

$\boxed{t(\rho)=\{(1,2),(1,4),(2,4),(3,1),(3,2),(3,4)\}}$.

(New pairs from $\rho^2$: $(1,4)$ via $1\!\to\!2\!\to\!4$; $(3,4)$ via $3\!\to\!2\!\to\!4$; $(3,2)$ via $3\!\to\!1\!\to\!2$ already present. Vertex 4 is a sink, so $\rho^3$ adds nothing.)

%% ────────────────────────────────────────

\subsection*{Problem 7.3}

\emph{Proof sketch.} $\tau=\rho\cap\sigma$ where $\rho,\sigma$ are equivalence relations.
Reflexive: $(x,x)\in\rho\cap\sigma$.
Symmetric: $(x,y)\in\tau\Rightarrow(y,x)\in\rho$ and $(y,x)\in\sigma\Rightarrow(y,x)\in\tau$.
Transitive: $(x,y),(y,z)\in\tau\subseteq\rho\Rightarrow(x,z)\in\rho$; similarly $\in\sigma$; so $(x,z)\in\tau$. \qed

%% ────────────────────────────────────────

\subsection*{Problem 7.4}

$\rho\circ\theta$ is an equivalence relation iff $\rho\circ\theta=\theta\circ\rho$.

\emph{Proof sketch.}
$(\Rightarrow)$ $(\rho\circ\theta)^{-1}=\theta^{-1}\circ\rho^{-1}=\theta\circ\rho$. Since $\rho\circ\theta$ is symmetric, $\rho\circ\theta=\theta\circ\rho$.

$(\Leftarrow)$ Reflexivity: $\Delta\subseteq\rho,\theta\Rightarrow\Delta\subseteq\rho\circ\theta$. Symmetry: $(\rho\circ\theta)^{-1}=\theta\circ\rho=\rho\circ\theta$. Transitivity: $(\rho\circ\theta)^2=\rho\circ(\theta\circ\rho)\circ\theta =\rho\circ(\rho\circ\theta)\circ\theta\subseteq\rho^2\circ\theta^2\subseteq\rho\circ\theta$. \qed

%% ────────────────────────────────────────

\subsection*{Problem 7.5}

$\rho\cup\theta$ is an equivalence relation iff $\rho\circ\theta=\theta\circ\rho$.

\emph{Proof sketch.}
$(\Rightarrow)$ If $\rho\cup\theta$ is an equivalence (hence transitive), then $\rho\circ\theta\subseteq(\rho\cup\theta)\circ(\rho\cup\theta)\subseteq\rho\cup\theta$. Similarly $\theta\circ\rho\subseteq\rho\cup\theta$. Since $\rho\cup\theta\subseteq\rho\circ\theta$ (as $\rho=\rho\circ\Delta\subseteq\rho\circ\theta$), we get $\rho\cup\theta=\rho\circ\theta$, and by symmetry also $=\theta\circ\rho$.

$(\Leftarrow)$ If $\rho\circ\theta=\theta\circ\rho$, then by Problem~7.4, $\rho\circ\theta$ is an equivalence relation. One shows $\rho\cup\theta=\rho\circ\theta$, which is then an equivalence relation. \qed
